% GENERATED FILE. Edit canonical lessons, then run npm run generate:books.
% canonical-source-hash: fnv1a64:8fc6380cb6707ad2
% canonical-lessons: GU-C235-bhasvu, GU-C235-karadvu, GU-C235-palvu, GU-C235-catvu, GU-C235-puchdi

\chapter{The dog that barked}
\label{ch:gu-the-dog-that-barked}
\hlchaptermodality{235}

\begin{chapteropening}
\textbf{By the end of this chapter:} \emph{I can say what a pet dog did.}

Complete the last lesson of chapter 235: I can say what a pet dog did.
\end{chapteropening}

\section[\texorpdfstring{\gu{ભસવું} (bhasvũ)}{ભસવું (bhasvũ)}]{\gu{ભસવું} (bhasvũ) — to bark}

\label{lesson:GU-C235-bhasvu}



\begin{quote}
Before the new one: say the Gujarati for sheep, then the Gujarati for to graze.
\end{quote}

\subsection*{You'll want to know: \gu{ભસવું}}
\textbf{\gu{ભસવું}} — \emph{bhasvũ} — \textquotedblleft{}to bark\textquotedblright{}.

\begin{quote}
\textbf{\gu{રાત્રે} \gu{કૂતરો} \gu{ભસ્યો}.} — \emph{rātre kūtro bhasyo} — \textquotedblleft{}The dog barked at night.\textquotedblright{}
\end{quote}

\begin{cousinweb}
\textbf{\gu{ભસવું}} (\emph{bhasvũ}) is a tadbhava, a word inherited from Sanskrit and reshaped by sound change. It comes from \emph{bhaṣati}, \textquotedblleft{}barks, growls\textquotedblright{} (CDIAL 9423).
\end{cousinweb}

\begin{grammarlens}[title={What you've built}]
One more word for everyday life at A2.
\end{grammarlens}

\subsection*{Your turn}
\begin{itemize}
\raggedright
  \item \emph{Say it:} \emph{bhasvũ}
  \item \emph{Say it:} \emph{bhasvũ}, once more
  \item \emph{Say it:} \emph{rātre kūtro bhasyo} — The dog barked at night.
  \item \emph{Recall:} say the Gujarati for sheep, then the Gujarati for to graze, then say \emph{bhasvũ} again
\end{itemize}

\begin{tcolorbox}[breakable,colback=teal!4,colframe=teal!35!black,title={Before you move on}]
What does \gu{ભસવું} mean? (\textquotedblleft{}To bark\textquotedblright{}.) Say it once more.
\end{tcolorbox}
\section[\texorpdfstring{\gu{કરડવું} (karaḍvũ)}{કરડવું (karaḍvũ)}]{\gu{કરડવું} (karaḍvũ) — to bite}

\label{lesson:GU-C235-karadvu}



\begin{quote}
Before the new one: say the Gujarati for to graze, then the Gujarati for to bark.
\end{quote}

\subsection*{You'll want to know: \gu{કરડવું}}
\textbf{\gu{કરડવું}} — \emph{karaḍvũ} — \textquotedblleft{}to bite\textquotedblright{}.

\begin{quote}
\textbf{\gu{ગઈકાલે} \gu{મને} \gu{કૂતરો} \gu{કરડ્યો}.} — \emph{gaīkāle mane kūtro karaḍyo} — \textquotedblleft{}Yesterday a dog bit me.\textquotedblright{}
\end{quote}

\begin{cousinweb}
The origin of \textbf{\gu{કરડવું}} (\emph{karaḍvũ}) is not settled: CDIAL lists no Gujarati form like it, and gives \emph{ḍasvũ}, \textquotedblleft{}bites, stings\textquotedblright{}, as the Gujarati word under Sanskrit \emph{daśati}, \textquotedblleft{}bites\textquotedblright{} (CDIAL 6230).
\end{cousinweb}

\begin{grammarlens}[title={What you've built}]
One more word for everyday life at A2.
\end{grammarlens}

\subsection*{Your turn}
\begin{itemize}
\raggedright
  \item \emph{Say it:} \emph{karaḍvũ}
  \item \emph{Say it:} \emph{karaḍvũ}, once more
  \item \emph{Say it:} \emph{gaīkāle mane kūtro karaḍyo} — Yesterday a dog bit me.
  \item \emph{Recall:} say the Gujarati for to graze, then the Gujarati for to bark, then say \emph{karaḍvũ} again
\end{itemize}

\begin{tcolorbox}[breakable,colback=teal!4,colframe=teal!35!black,title={Before you move on}]
What does \gu{કરડવું} mean? (\textquotedblleft{}To bite\textquotedblright{}.) Say it once more.
\end{tcolorbox}
\section[\texorpdfstring{\gu{પાળવું} (pāḷvũ)}{પાળવું (pāḷvũ)}]{\gu{પાળવું} (pāḷvũ) — to keep a pet, to rear; to obey}

\label{lesson:GU-C235-palvu}



\begin{quote}
Before the new one: say the Gujarati for to bark, then the Gujarati for to bite.
\end{quote}

\subsection*{You'll want to know: \gu{પાળવું}}
\textbf{\gu{પાળવું}} — \emph{pāḷvũ} — \textquotedblleft{}to keep a pet, to rear; to obey\textquotedblright{}.

\begin{quote}
\textbf{\gu{તેણે} \gu{એક} \gu{કૂતરો} \gu{પાળ્યો}.} — \emph{teṇe ek kūtro pāḷyo} — \textquotedblleft{}He kept a dog.\textquotedblright{}
\end{quote}

\begin{cousinweb}
\textbf{\gu{પાળવું}} (\emph{pāḷvũ}) is a tadbhava. It comes from Sanskrit \emph{pālayati}, \textquotedblleft{}guards, protects\textquotedblright{}, and CDIAL glosses the Gujarati verb \textquotedblleft{}to rear, keep\textquotedblright{} (CDIAL 8129).
\end{cousinweb}

\begin{grammarlens}[title={What you've built}]
One more word for everyday life at A2.
\end{grammarlens}

\subsection*{Your turn}
\begin{itemize}
\raggedright
  \item \emph{Say it:} \emph{pāḷvũ}
  \item \emph{Say it:} \emph{pāḷvũ}, once more
  \item \emph{Say it:} \emph{teṇe ek kūtro pāḷyo} — He kept a dog.
  \item \emph{Recall:} say the Gujarati for to bark, then the Gujarati for to bite, then say \emph{pāḷvũ} again
\end{itemize}

\begin{tcolorbox}[breakable,colback=teal!4,colframe=teal!35!black,title={Before you move on}]
What does \gu{પાળવું} mean? (\textquotedblleft{}To keep a pet, to rear; to obey\textquotedblright{}.) Say it once more.
\end{tcolorbox}
\section[\texorpdfstring{\gu{ચાટવું} (cāṭvũ)}{ચાટવું (cāṭvũ)}]{\gu{ચાટવું} (cāṭvũ) — to lick}

\label{lesson:GU-C235-catvu}



\begin{quote}
Before the new one: say the Gujarati for to bite, then the Gujarati for to keep a pet.
\end{quote}

\subsection*{You'll want to know: \gu{ચાટવું}}
\textbf{\gu{ચાટવું}} — \emph{cāṭvũ} — \textquotedblleft{}to lick\textquotedblright{}.

\begin{quote}
\textbf{\gu{તેણે} \gu{મારો} \gu{હાથ} \gu{ચાટ્યો}.} — \emph{teṇe māro hāth cāṭyo} — \textquotedblleft{}It licked my hand.\textquotedblright{}
\end{quote}

\begin{cousinweb}
\textbf{\gu{ચાટવું}} (\emph{cāṭvũ}) comes from \emph{*caṭṭ}, \textquotedblleft{}lick, taste\textquotedblright{}, a reconstructed form that CDIAL thinks may be imitative; it lists the Gujarati verb, glossed \textquotedblleft{}to lick\textquotedblright{} (CDIAL 4573).
\end{cousinweb}

\begin{grammarlens}[title={What you've built}]
One more word for everyday life at A2.
\end{grammarlens}

\subsection*{Your turn}
\begin{itemize}
\raggedright
  \item \emph{Say it:} \emph{cāṭvũ}
  \item \emph{Say it:} \emph{cāṭvũ}, once more
  \item \emph{Say it:} \emph{teṇe māro hāth cāṭyo} — It licked my hand.
  \item \emph{Recall:} say the Gujarati for to bite, then the Gujarati for to keep a pet, then say \emph{cāṭvũ} again
\end{itemize}

\begin{tcolorbox}[breakable,colback=teal!4,colframe=teal!35!black,title={Before you move on}]
What does \gu{ચાટવું} mean? (\textquotedblleft{}To lick\textquotedblright{}.) Say it once more.
\end{tcolorbox}
\section[\texorpdfstring{\gu{પૂંછડી} (pū̃chḍī)}{પૂંછડી (pū̃chḍī)}]{\gu{પૂંછડી} (pū̃chḍī) — tail}

\label{lesson:GU-C235-puchdi}



\begin{quote}
Before the new one: say the Gujarati for to keep a pet, then the Gujarati for to lick.
\end{quote}

\subsection*{You'll want to know: \gu{પૂંછડી}}
\textbf{\gu{પૂંછડી}} — \emph{pū̃chḍī} — \textquotedblleft{}tail\textquotedblright{}.

\begin{quote}
\textbf{\gu{કૂતરાની} \gu{પૂંછડી} \gu{લાંબી} \gu{હતી}.} — \emph{kūtrānī pū̃chḍī lā̃bī hatī} — \textquotedblleft{}The dog's tail was long.\textquotedblright{}
\end{quote}

\begin{cousinweb}
\textbf{\gu{પૂંછડી}} (\emph{pū̃chḍī}) is a tadbhava. It comes from Sanskrit \emph{puccha}, \textquotedblleft{}tail\textquotedblright{}, where CDIAL lists Gujarati \emph{pũch} and \emph{puchṛũ}, \textquotedblleft{}tail\textquotedblright{} (CDIAL 8249).
\end{cousinweb}

\begin{grammarlens}[title={What you've built}]
One more word for everyday life at A2.
\end{grammarlens}

\subsection*{Your turn}
\begin{itemize}
\raggedright
  \item \emph{Say it:} \emph{pū̃chḍī}
  \item \emph{Say it:} \emph{pū̃chḍī}, once more
  \item \emph{Say it:} \emph{kūtrānī pū̃chḍī lā̃bī hatī} — The dog's tail was long.
  \item \emph{Recall:} say the Gujarati for to keep a pet, then the Gujarati for to lick, then say \emph{pū̃chḍī} again
\end{itemize}

\begin{tcolorbox}[breakable,colback=teal!4,colframe=teal!35!black,title={Before you move on}]
What does \gu{પૂંછડી} mean? (\textquotedblleft{}Tail\textquotedblright{}.) Say it once more.
\end{tcolorbox}
